\documentclass[noauthor]{ximera}

\title{Week 12}
\author{Math 1193 Design Team}

\begin{document}
\begin{abstract}
    Activities for Week 12.
\end{abstract}
\maketitle

\begin{problem}
Which of the following expressions are equal to \(abc\)? There are multiple
answers.
\begin{enumerate}
    \item \(\left((abc)^{-1}\right)^{-1}\)
    \item \((abc)^{-1}a^{-1}b^{-1}c^{-1}\)
    \item \(\frac{bc}{\frac{1}{a^{-1}}}\)
    \item \((abc + 1)^2 - 1\)
    \item \(\frac{\sqrt{(ab)^2}}{c^{-2}}\)
    \item \(\sqrt{a}\sqrt{b}\sqrt{c}a^{1/2}b^{1/2}c^{1/2}\)
    \item \(\frac{(a^{2}b^{2})^{1/2}}{c^{-1}}\)
    \item \(\frac{\sqrt{(-a)^2}((cb)^2 - 1)}{cb + 1} + a\)
    \item \(\left(\frac{a^2b^2}{c^{-2}}\right)^{1/2}\)
\end{enumerate}
\begin{instructorNotes}
    This is meant to get around students just plugging the expressions into their calculators. 
\end{instructorNotes}
\end{problem}

\begin{problem}
Let \(f\) be defined by \(f(x) = \frac{3x}{x + 1}\).

\begin{enumerate}
\item \(f\) is invertible. How can you tell?
\item Is \(g(x) = \frac{3x}{x^2 + 1}\) invertible?
\item Find the domain of \(f\).
\item Evaluate the following expressions: 
\begin{enumerate}
    \item \(f^{-1}(f(2))\)
    \item \(f^{-1}(f(-1))\)
\end{enumerate}
\item Find a formula for \(f^{-1}(x)\).
\item Find the range of \(f\).
\end{enumerate}

\begin{instructorNotes}
    Students have issues with inverse functions and the relationship between domain and range of inverses.
    This is also a difficult one for students to find the inverse of, since there are multiple instances of \(x\).
\end{instructorNotes}
\end{problem}



\begin{problem}
\includegraphics[width=10cm]{base3.png}
\begin{enumerate}
    \item Above is the graph of \(g(x) = 3^x\). Label at least
    three points on the graph.
    \item Reflect your graph across the line \(y = x\). To do this, swap the \(x\) and \(y\) coordinates of each
    point on the graph. Call your new function \(f\).
    \begin{enumerate}
        \item What is the relationship between \(g\) and \(f\)?
        \item What are three points on the graph of \(f\)?
        \item What is another name for \(f\)? There are two possible answers:
        see if you can come up with both!
    \end{enumerate}
    \item Is the asympote to the graph of \(f\) vertical or horizontal?
    What is its equation?
    \item What are the domain and range of \(f\)? If you get
    stuck, review the relationship between the domains and 
    ranges of inverse functions.
    
\end{enumerate}

\begin{instructorNotes}
This is meant to draw on knowledge of inverse functions and also practice with the graphs of exponentials.
\end{instructorNotes}
\end{problem}

\begin{problem}
    \begin{enumerate}
        \item Help! I made a table containing some points on the graph of \(f(x) = x^2\), but my inconsiderate roommate erased parts of the table. Please fill it in for me. My professor is very picky about exact answers, so don't round! Are there any entries that could have more than one option?
        \begin{center}
\begin{tabular}{|c|c|}
 \(x\) & \(y\) \\ \hline 
 1 &  \\  
 \(-2\) & 4 \\
  & 4 \\
  & 2 \\
 0 & \\
 \(-9\) & \\
 & 3 \\
 \(\sqrt{2}\) & \\
\end{tabular}
\end{center}
        \item Same thing, but for \(g(x) = e^x\). 
                \begin{center}
\begin{tabular}{|c|c|}
 \(x\) & \(y\) \\ \hline 
 1 &  \\  
 2 & \(e^2\) \\
  & \(e\) \\
 0 & \\
 \(-9\) & \\
 & \(\sqrt{e}\) \\
  & \(\frac{1}{e^2}\) \\
  \(-\frac{1}{3}\) & \\
\end{tabular}
\end{center}
    \end{enumerate}
\begin{instructorNotes}
    This is meant to convince students to be comfortable with square roots of non-square numbers and powers of \(e\).
\end{instructorNotes}
\end{problem}

\begin{problem}
Aaliyah and Tom are having a friendly disagreement. Aaliyah thinks
\(\log(15) = \log(3) + \log(5)\) and Tom thinks
\(\log(15) = \log(3) \cdot \log(5)\).

Aaliyah's work is as follows:
    \begin{align*}
        \log(15) & = \log(3 \cdot 5) \\
        & = \log(10^{\log(3)} \cdot 10^{\log(5)}) \\
        & = \log(10^{\log(3) + \log(5)}) \\
        & = \log(3) + \log(5)
    \end{align*}

Tom's work is as follows:
    \begin{align*}
        \log(15) & = \log(3 \cdot 5) \\
        & = \log(10^{\log(3)} \cdot 10^{\log(5)}) \\
        & = \log(10^{\log(3) \cdot \log(5)}) \\
        & = \log(3) \cdot \log(5)
    \end{align*}


Whose work is correct? Why? How could you check?

\end{problem}

\begin{problem}
When I was learning logarithms, there was no way to put \(\log_2(5)\)
into my calculator, since there was only a \(\log\) button and a
\(\ln\) button, without the option to directly input a base other
than 10 or \(e\). In this exercise, we'll learn how my generation
dealt with this problem. 

Make sure to write down your work at each step.

\begin{enumerate}
    \item Write the equation \(\log_2(5) = x\) in exponential form.
    \item Now take the \textbf{natural log} of both sides of your
    exponential equation.
    \item Next, we'll use a property of logarithms we'll learn soon:
    \(\ln(2^x) = x \ln(2)\). Use this to rewrite one side of your equation.
    \item Finally, solve for \(x\).
    \item How did I put \(\log_2(5)\) into my calculator? What about
    \(\log_7(19)\)?
\end{enumerate}
\end{problem}

\begin{problem}
\begin{enumerate}
    \item Sky was asked to simplify \(\ln(e^{-2})\) and said the answer was
    \(-2\). Are they correct? Why or why not? Write down your reasoning.
    \item Sky was asked to simplify \(e^{\ln(-2)}\) and said the answer was
    \(-2\). Are they correct? Why or why not? Write down your reasoning.
\end{enumerate}
\end{problem}

\end{document}
